% GENERATED FILE. Edit canonical lessons, then run npm run generate:books.
% canonical-source-hash: fnv1a64:1bba42a4c335d16e
% canonical-lessons: FR-C509-plonger, FR-C509-plongeon, FR-C509-maillot-de-bain, FR-C509-brasse, FR-C509-ramer

\chapter{Dans l'eau}
\label{ch:fr-dans-l-eau}
\hlchaptermodality{509}

\begin{chapteropening}
\textbf{By the end of this chapter:} \emph{I can talk about swimming, diving and rowing.}

Complete the last lesson of chapter 509: I can talk about swimming, diving and rowing.
\end{chapteropening}

\section[plonger]{plonger — to dive}

\label{lesson:FR-C509-plonger}



\begin{quote}
Before the new one: say the French for relay, then the French for to time.
\end{quote}

\subsection*{You'll want to know: plonger}
\textbf{plonger} — \textquotedblleft{}to dive\textquotedblright{}.

\begin{quote}
Les enfants plongent du haut du rocher.
\end{quote}

\begin{quote}
\emph{The children dive off the top of the rock.}
\end{quote}

Spelling: \emph{nous plongeons}. \emph{La plongée} — scuba diving.

\begin{cousinweb}
\emph{plonger} is inherited from Old French \emph{plongier} (EtymDB), from Vulgar Latin \emph{plumbicāre} \textquotedblleft{}to sink like lead\textquotedblright{}, built on \emph{plumbum} \textquotedblleft{}lead\textquotedblright{} (Lewis \& Short).
\end{cousinweb}

\begin{grammarlens}[title={What you've built}]
One more word for telling, explaining and describing at B1.
\end{grammarlens}

\subsection*{Your turn}
\begin{itemize}
\raggedright
  \item \emph{Say it:} \emph{plonger}
  \item \emph{Say it:} \emph{plonger}, once more
  \item \emph{Say it:} a sentence of your own with \emph{plonger}
  \item \emph{Recall:} say the French for relay, then the French for to time, then say \emph{plonger} again
\end{itemize}

\begin{tcolorbox}[breakable,colback=teal!4,colframe=teal!35!black,title={Before you move on}]
What does plonger mean? (\textquotedblleft{}To dive\textquotedblright{}.) Say it once more.
\end{tcolorbox}
\section[le plongeon]{le plongeon — dive}

\label{lesson:FR-C509-plongeon}



\begin{quote}
Before the new one: say the French for to time, then the French for to dive.
\end{quote}

\subsection*{You'll want to know: le plongeon}
\textbf{le plongeon} — \textquotedblleft{}dive\textquotedblright{}.

\begin{quote}
Il a fait un beau plongeon dans la piscine.
\end{quote}

\begin{quote}
\emph{He did a nice dive into the pool.}
\end{quote}

\emph{Faire un plongeon} — to dive in.

\begin{cousinweb}
\emph{le plongeon} is built on \emph{plonger} \textquotedblleft{}to dive\textquotedblright{} (EtymDB), from Vulgar Latin \emph{plumbicāre} \textquotedblleft{}to sink like lead\textquotedblright{}, built on \emph{plumbum} \textquotedblleft{}lead\textquotedblright{} (Lewis \& Short).
\end{cousinweb}

\begin{grammarlens}[title={What you've built}]
One more word for telling, explaining and describing at B1.
\end{grammarlens}

\subsection*{Your turn}
\begin{itemize}
\raggedright
  \item \emph{Say it:} \emph{le plongeon}
  \item \emph{Say it:} \emph{le plongeon}, once more
  \item \emph{Say it:} a sentence of your own with \emph{le plongeon}
  \item \emph{Recall:} say the French for to time, then the French for to dive, then say \emph{le plongeon} again
\end{itemize}

\begin{tcolorbox}[breakable,colback=teal!4,colframe=teal!35!black,title={Before you move on}]
What does le plongeon mean? (\textquotedblleft{}Dive\textquotedblright{}.) Say it once more.
\end{tcolorbox}
\section[le maillot de bain]{le maillot de bain — swimsuit, swimming trunks}

\label{lesson:FR-C509-maillot-de-bain}



\begin{quote}
Before the new one: say the French for to dive, then the French for dive.
\end{quote}

\subsection*{You'll want to know: le maillot de bain}
\textbf{le maillot de bain} — \textquotedblleft{}swimsuit, swimming trunks\textquotedblright{}.

\begin{quote}
J'ai oublié mon maillot de bain à l'hôtel.
\end{quote}

\begin{quote}
\emph{I left my swimsuit at the hotel.}
\end{quote}

Often \emph{le maillot} alone.

\begin{cousinweb}
\emph{le maillot de bain} joins \emph{maillot}, built on \emph{maille} \textquotedblleft{}a stitch, a mesh\textquotedblright{}, from Latin \emph{macula} \textquotedblleft{}a spot, a mesh\textquotedblright{}, and \emph{bain} \textquotedblleft{}bath\textquotedblright{}, from Latin \emph{balneum} (Lewis \& Short).
\end{cousinweb}

\begin{grammarlens}[title={What you've built}]
One more word for telling, explaining and describing at B1.
\end{grammarlens}

\subsection*{Your turn}
\begin{itemize}
\raggedright
  \item \emph{Say it:} \emph{le maillot de bain}
  \item \emph{Say it:} \emph{le maillot de bain}, once more
  \item \emph{Say it:} a sentence of your own with \emph{le maillot de bain}
  \item \emph{Recall:} say the French for to dive, then the French for dive, then say \emph{le maillot de bain} again
\end{itemize}

\begin{tcolorbox}[breakable,colback=teal!4,colframe=teal!35!black,title={Before you move on}]
What does le maillot de bain mean? (\textquotedblleft{}Swimsuit, swimming trunks\textquotedblright{}.) Say it once more.
\end{tcolorbox}
\section[la brasse]{la brasse — breaststroke}

\label{lesson:FR-C509-brasse}



\begin{quote}
Before the new one: say the French for dive, then the French for swimsuit.
\end{quote}

\subsection*{You'll want to know: la brasse}
\textbf{la brasse} — \textquotedblleft{}breaststroke\textquotedblright{}.

\begin{quote}
Ma grand-mère nage la brasse tous les matins.
\end{quote}

\begin{quote}
\emph{My grandmother swims breaststroke every morning.}
\end{quote}

Front crawl is \emph{le crawl}; backstroke is \emph{le dos crawlé}.

\begin{cousinweb}
\emph{la brasse} is inherited from Latin \emph{bracchia} \textquotedblleft{}arms\textquotedblright{} (EtymDB), the plural of \emph{bracchium} \textquotedblleft{}arm\textquotedblright{} (Lewis \& Short): first the span of two outstretched arms.
\end{cousinweb}

\begin{grammarlens}[title={What you've built}]
One more word for telling, explaining and describing at B1.
\end{grammarlens}

\subsection*{Your turn}
\begin{itemize}
\raggedright
  \item \emph{Say it:} \emph{la brasse}
  \item \emph{Say it:} \emph{la brasse}, once more
  \item \emph{Say it:} a sentence of your own with \emph{la brasse}
  \item \emph{Recall:} say the French for dive, then the French for swimsuit, then say \emph{la brasse} again
\end{itemize}

\begin{tcolorbox}[breakable,colback=teal!4,colframe=teal!35!black,title={Before you move on}]
What does la brasse mean? (\textquotedblleft{}Breaststroke\textquotedblright{}.) Say it once more.
\end{tcolorbox}
\section[ramer]{ramer — to row}

\label{lesson:FR-C509-ramer}



\begin{quote}
Before the new one: say the French for swimsuit, then the French for breaststroke.
\end{quote}

\subsection*{You'll want to know: ramer}
\textbf{ramer} — \textquotedblleft{}to row\textquotedblright{}.

\begin{quote}
Nous avons ramé jusqu'à l'île.
\end{quote}

\begin{quote}
\emph{We rowed to the island.}
\end{quote}

Informally, to struggle: \emph{j'ai ramé pour trouver un appartement}.

\begin{cousinweb}
\emph{ramer} is linked to \emph{rame} \textquotedblleft{}oar\textquotedblright{}. EtymDB traces it to Old French \emph{raim} \textquotedblleft{}branch\textquotedblright{}, from Latin \emph{rāmus} \textquotedblleft{}a branch\textquotedblright{}; Latin \emph{rēmus} \textquotedblleft{}an oar\textquotedblright{} (Lewis \& Short) probably shaped the sense, and the origin is not fully settled.
\end{cousinweb}

\begin{grammarlens}[title={What you've built}]
One more word for telling, explaining and describing at B1.
\end{grammarlens}

\subsection*{Your turn}
\begin{itemize}
\raggedright
  \item \emph{Say it:} \emph{ramer}
  \item \emph{Say it:} \emph{ramer}, once more
  \item \emph{Say it:} a sentence of your own with \emph{ramer}
  \item \emph{Recall:} say the French for swimsuit, then the French for breaststroke, then say \emph{ramer} again
\end{itemize}

\begin{tcolorbox}[breakable,colback=teal!4,colframe=teal!35!black,title={Before you move on}]
What does ramer mean? (\textquotedblleft{}To row\textquotedblright{}.) Say it once more.
\end{tcolorbox}
